\section{Conclusion}

\subsection{Conclusion}

This study developed a multiple linear regression model to examine factors associated with student academic performance. Using the 60-observation synthetic dataset, the model considered study time, attendance, previous performance, and access to learning resources as explanatory variables.

The fitted model showed positive relationships between all four explanatory variables and academic performance. Study time, attendance, previous performance, and learning resources were all statistically significant at the 5\% level. The model explained approximately 78.2\% of the variation in academic performance, while the adjusted $R^2$ was 76.6\%.

Overall, the results demonstrate how multiple linear regression can be used to examine several factors simultaneously and provide a framework for understanding and predicting student academic performance.

\subsection{Recommendations}

Based on the findings of the study, the following recommendations are made:

\begin{enumerate}
    \item Students should be encouraged to develop consistent and effective study habits.
    \item Educational institutions should promote regular attendance and identify students with attendance difficulties.
    \item Previous academic performance should be considered when identifying students who may require additional academic support.
    \item Students should be provided with adequate learning resources to support their academic activities.
    \item Statistical models can be used as supplementary tools to help educators identify factors associated with student performance.
\end{enumerate}

\subsection{Limitations}

The dataset used in this study is synthetic and was created for demonstrative purposes. Therefore, the findings should not be interpreted as evidence about a real student population.

In addition, the model considers only the variables included in the dataset. Other factors that may influence academic performance were not examined. The regression results should therefore be viewed as an illustration of the statistical modelling process rather than as definitive conclusions about student academic performance.